% ================================================================
% DEEP LEARNING FUNDAMENTALS
% Comprehensive introduction to neural networks and deep learning
%
% Topics Covered:
% 1. Neural Network Foundations
% 2. Training Deep Networks (Backprop, Optimization)
% 3. Convolutional Neural Networks (CNNs)
% 4. Recurrent Neural Networks (RNNs)
% 5. Modern Architectures (Transformers, ResNets)
% 6. Practical Considerations
%
% Total: ~45 slides
% Author: Diogo Ribeiro
% Date: January 2025
% ================================================================

\section{Deep Learning Fundamentals}

% ================================================================
% PART 1: NEURAL NETWORK BASICS
% ================================================================

\begin{frame}{From Linear Models to Neural Networks}
\textbf{Linear regression:}
\[
  \hat{y} = w_1 x_1 + w_2 x_2 + \cdots + w_d x_d + b
\]

\begin{itemize}
  \item Linear models represent additive linear relationships in the chosen features.
  \item Non-linear activations let multilayer networks represent non-linear decision surfaces.
\end{itemize}

\textbf{Adding non-linearity:}
\[
  \hat{y} = \sigma(w_1 x_1 + w_2 x_2 + b).
\]

\begin{alertblock}{Universal approximation}
Under standard universal-approximation conditions, a sufficiently wide network with an appropriate non-linear activation can approximate continuous functions on compact sets arbitrarily well. This does not imply efficient learning, finite-sample generalization, or practical trainability.
\end{alertblock}
\end{frame}

\begin{frame}{From a Single Neuron to an MLP}
\begin{columns}
\begin{column}{0.48\textwidth}
\textbf{Single neuron}
\begin{figure}
\centering
\begin{tikzpicture}[scale=0.75]
\node[draw, circle] (x1) at (0,2) {$x_1$};
\node[draw, circle] (x2) at (0,1) {$x_2$};
\node[draw, circle] (x3) at (0,0) {$x_3$};
\node[draw, circle, fill=navyblue!20, minimum size=1.4cm] (n) at (3,1) {$\sigma$};
\node[draw, circle] (y) at (5.5,1) {$\hat{y}$};
\draw[->, thick] (x1) -- node[above] {$w_1$} (n);
\draw[->, thick] (x2) -- node[above] {$w_2$} (n);
\draw[->, thick] (x3) -- node[above] {$w_3$} (n);
\draw[->, thick] (n) -- (y);
\end{tikzpicture}
\end{figure}
\end{column}

\begin{column}{0.48\textwidth}
\textbf{Multi-layer perceptron}
\begin{figure}
\centering
\begin{tikzpicture}[scale=0.75]
\foreach \i in {1,...,3} {
  \node[draw, circle, fill=forest!20] (x\i) at (0,{3-\i}) {$x_{\i}$};
}
\foreach \i in {1,...,4} {
  \node[draw, circle, fill=navyblue!20] (h\i) at (2.5,{3.5-\i}) {};
}
\node[draw, circle, fill=crimson!20] (y) at (5,1) {$\hat{y}$};
\foreach \i in {1,...,3} {
  \foreach \j in {1,...,4} {
    \draw[->] (x\i) -- (h\j);
  }
}
\foreach \i in {1,...,4} { \draw[->] (h\i) -- (y); }
\node at (0,-1) {Input};
\node at (2.5,-1) {Hidden};
\node at (5,-1) {Output};
\end{tikzpicture}
\end{figure}
\end{column}
\end{columns}

\begin{block}{Depth}
Deep learning composes multiple representation-learning layers; depth is not the same thing as parameter count.
\end{block}
\end{frame}

\begin{frame}{Activation Functions: Saturating Nonlinearities}
Without a non-linearity, stacked affine layers collapse to one affine transformation.

\textbf{Sigmoid}
\[
  \sigma(z)=\frac{1}{1+e^{-z}}
\]
\begin{itemize}
  \item useful for Bernoulli output probabilities;
  \item saturates for large $|z|$, producing small derivatives.
\end{itemize}

\textbf{Tanh}
\begin{itemize}
  \item output in $(-1,1)$ and centered around zero;
  \item also saturates for large $|z|$.
\end{itemize}
\end{frame}

\begin{frame}{Activation Functions: ReLU Family}
\textbf{ReLU}
\[
  \operatorname{ReLU}(z)=\max(0,z).
\]
\begin{itemize}
  \item simple and non-saturating for positive inputs;
  \item has zero gradient for negative inputs and can create inactive units;
  \item does not eliminate vanishing-gradient problems across an entire network.
\end{itemize}

\textbf{Variants}
\begin{itemize}
  \item Leaky ReLU and ELU modify the negative branch;
  \item GELU is common in Transformer architectures;
  \item SiLU/Swish uses $z\sigma(z)$.
\end{itemize}

\begin{alertblock}{Choice}
Activation choice interacts with architecture, initialization, normalization, and optimization. There is no universal default for every network.
\end{alertblock}
\end{frame}

\begin{frame}{Backpropagation: The Optimization Problem}
Given training data $\{(x^{(i)},y^{(i)})\}_{i=1}^n$, minimize
\[
  \mathcal{L}(W)=\frac{1}{n}\sum_{i=1}^n \operatorname{loss}(\hat y^{(i)},y^{(i)}).
\]

Gradient-based optimization updates parameters as
\[
  W \leftarrow W - \eta\nabla_W\mathcal{L}.
\]

For a two-layer network
\[
  \hat y=f\!\left(W_2\,\sigma(W_1x)\right),
\]
backpropagation applies the chain rule efficiently through the computational graph.

\begin{block}{Key distinction}
Backpropagation computes derivatives. The optimizer decides how those derivatives update parameters.
\end{block}
\end{frame}

\begin{frame}{Backpropagation: Forward and Backward Passes}
\begin{columns}
\begin{column}{0.42\textwidth}
\textbf{Forward pass}
\begin{enumerate}
  \item compute layer activations;
  \item store intermediate quantities;
  \item evaluate the loss.
\end{enumerate}

\textbf{Backward pass}
\begin{enumerate}
  \item propagate sensitivities from the loss;
  \item accumulate parameter gradients;
  \item pass gradients to the optimizer.
\end{enumerate}
\end{column}

\begin{column}{0.56\textwidth}
\begin{figure}
\centering
\begin{tikzpicture}[scale=0.72]
\node[draw, rectangle, fill=forest!20] (x) at (0,3) {$x$};
\node[draw, rectangle, fill=navyblue!20] (h) at (2,3) {$h$};
\node[draw, rectangle, fill=crimson!20] (y) at (4,3) {$\hat y$};
\node[draw, rectangle, fill=purple!20] (l) at (6,3) {$\mathcal L$};
\draw[->, thick] (x) -- node[above] {$W_1$} (h);
\draw[->, thick] (h) -- node[above] {$W_2$} (y);
\draw[->, thick] (y) -- (l);
\node at (3,3.7) {Forward};
\draw[<-, thick, crimson] (0,1) -- node[below] {$\partial\mathcal L/\partial W_1$} (2,1);
\draw[<-, thick, crimson] (2,1) -- node[below] {$\partial\mathcal L/\partial W_2$} (4,1);
\draw[<-, thick, crimson] (4,1) -- node[below] {$\partial\mathcal L/\partial\hat y$} (6,1);
\node at (3,0.3) {Backward};
\end{tikzpicture}
\end{figure}
\end{column}
\end{columns}

\begin{alertblock}{Automatic differentiation}
Modern frameworks record the computational graph and evaluate the required derivatives automatically.
\end{alertblock}
\end{frame}

\begin{frame}{Optimization: SGD and Momentum}
\begin{itemize}
  \item Full-batch gradient descent uses the complete training set per update; minibatch SGD uses a stochastic gradient estimate.
  \item Minibatches reduce memory cost and increase update frequency, but introduce gradient noise whose effect depends on batch size, learning rate, and objective geometry.
\end{itemize}

\textbf{Momentum:}
\[
  v_t=\gamma v_{t-1}+\eta g_t,
  \qquad W_t=W_{t-1}-v_t.
\]

\begin{block}{Why momentum helps}
Momentum can smooth oscillatory directions and accelerate progress along persistent gradient directions. It does not guarantee escape from poor local geometry.
\end{block}
\end{frame}

\begin{frame}[fragile]{Adaptive Optimizers and Learning-Rate Schedules}
\begin{columns}
\begin{column}{0.47\textwidth}
\textbf{Common choices}
\begin{itemize}
  \item Adam: adaptive per-parameter scaling plus momentum-like moments;
  \item AdamW: decouples weight decay from the adaptive update;
  \item SGD + momentum: often competitive with a tuned schedule.
\end{itemize}

\textbf{Schedules}
\begin{itemize}
  \item step or exponential decay;
  \item cosine schedules;
  \item warm-up or one-cycle policies where justified.
\end{itemize}
\end{column}

\begin{column}{0.49\textwidth}
\begin{lstlisting}[style=code,basicstyle=\ttfamily\small]
optimizer = torch.optim.AdamW(
    model.parameters(),
    lr=3e-4,
    weight_decay=1e-2,
)

scheduler = torch.optim.lr_scheduler.
    CosineAnnealingLR(
        optimizer, T_max=100
    )
\end{lstlisting}
\end{column}
\end{columns}

\begin{alertblock}{No universal default}
Optimizer, learning rate, schedule, batch size, and regularization interact. Compare plausible configurations under matched training budgets.
\end{alertblock}
\end{frame}

% ================================================================
% PART 2: CNNs
% ================================================================

\begin{frame}{CNNs: Local Connectivity and Parameter Sharing}
Convolutional layers exploit spatial locality by applying the same learned kernel across positions.

\begin{itemize}
  \item local receptive fields reduce the number of free parameters relative to dense connections;
  \item weight sharing makes convolution translation-equivariant: shifting the input shifts the feature map;
  \item learned filters can detect increasingly complex spatial patterns across depth.
\end{itemize}

\[
  (X*K)_{ij}=\sum_m\sum_n X_{i+m,j+n}K_{mn}.
\]

\begin{alertblock}{Invariance}
Convolution itself is equivariant, not invariant. Pooling, striding, augmentation, and architecture choices can create approximate invariance to selected transformations.
\end{alertblock}
\end{frame}

\begin{frame}{CNNs: From Feature Maps to Prediction}
\begin{itemize}
  \item convolution produces feature maps;
  \item nonlinearities transform those maps;
  \item pooling or strided convolutions reduce spatial resolution;
  \item deeper layers increase the effective receptive field;
  \item a prediction head maps the learned representation to the task output.
\end{itemize}

\[
  \text{Input}\rightarrow
  [\text{Conv}\rightarrow\text{Activation}]^L
  \rightarrow\text{Head}\rightarrow\text{Output}.
\]

\begin{block}{Modern practice}
Many current CNNs use residual connections, normalization, global pooling, and task-specific heads rather than a fixed Conv--ReLU--Pool recipe.
\end{block}
\end{frame}

\begin{frame}[fragile,allowframebreaks]{Modern CNN Architectures}
\textbf{Evolution from LeNet to EfficientNet}

\begin{columns}
\begin{column}{0.5\textwidth}
\textbf{Classic Architectures:}

\begin{enumerate}
\item \textbf{LeNet-5 (1998):}
  \begin{itemize}
  \item 7 layers, digit recognition
  \item First successful CNN
  \end{itemize}

\item \textbf{AlexNet (2012):}
  \begin{itemize}
  \item ImageNet winner
  \item 8 layers, 60M parameters
  \item ReLU, dropout, GPU training
  \end{itemize}

\item \textbf{VGG (2014):}
  \begin{itemize}
  \item 16-19 layers
  \item Small 3×3 filters stacked
  \item Very deep, simple architecture
  \end{itemize}

\item \textbf{ResNet (2015):}
  \begin{itemize}
  \item 50-152 layers
  \item \textcolor{crimson}{Skip connections!}
  \item Solves vanishing gradient
  \end{itemize}
\end{enumerate}

\textbf{ResNet Skip Connection:}
\[
y = F(x, W) + x
\]

Identity mapping helps gradient flow!
\end{column}

\begin{column}{0.5\textwidth}
\textbf{Modern Developments:}

\begin{table}
\tiny
\begin{tabular}{|l|p{4.5cm}|}
\hline
\rowcolor{navyblue!20}
\textbf{Architecture} & \textbf{Innovation} \\
\hline
\textbf{Inception} & Multiple filter sizes in parallel \\
\hline
\textbf{ResNet} & Skip connections (residual blocks) \\
\hline
\textbf{DenseNet} & Every layer connects to every other \\
\hline
\textbf{MobileNet} & Depthwise separable convs (mobile) \\
\hline
\textbf{EfficientNet} & Compound scaling (width+depth+resolution) \\
\hline
\textbf{Vision Transformer} & Self-attention instead of convolutions \\
\hline
\end{tabular}
\end{table}

\vspace{0.3cm}

\textbf{Transfer Learning:}

Instead of training from scratch:
\begin{enumerate}
\item Load pre-trained model (ImageNet)
\item Remove final layer
\item Add new layer for your task
\item Fine-tune on your data
\end{enumerate}
\end{column}
\end{columns}
\end{frame}

% ================================================================
% PART 3: RNNs and Transformers
% ================================================================

\begin{frame}{Recurrent Neural Networks (RNNs)}
\textbf{Networks with memory for sequential data}

\begin{columns}
\begin{column}{0.5\textwidth}
\textbf{Why RNNs?}

For sequences (text, time series, speech):
\begin{itemize}
\item Input length varies
\item Temporal dependencies matter
\item Order is important
\end{itemize}

\textbf{RNN Idea:}

Maintain hidden state $h_t$ that captures history:
\[
h_t = \tanh(W_{hh} h_{t-1} + W_{xh} x_t)
\]
\[
y_t = W_{hy} h_t
\]

Hidden state $h_t$ is updated at each timestep.

\vspace{0.3cm}

\textbf{Problem: Vanishing Gradients}

During backprop through time, gradients decay exponentially $\to$ can't learn long-term dependencies.
\end{column}

\begin{column}{0.5\textwidth}
\textbf{LSTM (Long Short-Term Memory):}

Solution to vanishing gradients:

\begin{itemize}
\item \textbf{Cell state:} Highway for information flow
\item \textbf{Gates:} Control info flow
  \begin{itemize}
  \item Forget gate: What to forget
  \item Input gate: What to update
  \item Output gate: What to output
  \end{itemize}
\end{itemize}

\vspace{0.2cm}

\textbf{LSTM Cell:}
\begin{align*}
f_t &= \sigma(W_f [h_{t-1}, x_t]) \quad \text{(forget)} \\
i_t &= \sigma(W_i [h_{t-1}, x_t]) \quad \text{(input)} \\
\tilde{C}_t &= \tanh(W_C [h_{t-1}, x_t]) \\
C_t &= f_t \odot C_{t-1} + i_t \odot \tilde{C}_t \\
o_t &= \sigma(W_o [h_{t-1}, x_t]) \quad \text{(output)} \\
h_t &= o_t \odot \tanh(C_t)
\end{align*}

\vspace{0.2cm}

\textbf{GRU (Gated Recurrent Unit):}

Simpler than LSTM, fewer parameters, often works as well.
\end{column}
\end{columns}
\end{frame}

\begin{frame}{Transformers: Self-Attention}
Transformers replace recurrent state updates with attention over a set or sequence of representations.

\[
  \operatorname{Attention}(Q,K,V)
  = \operatorname{softmax}\!\left(
      \frac{QK^\top}{\sqrt{d_k}}
    \right)V.
\]

\begin{itemize}
  \item queries specify what information is requested;
  \item keys determine compatibility with each query;
  \item values supply the information being aggregated;
  \item positional information is required when order matters.
\end{itemize}

\begin{block}{Trade-off}
Self-attention supports parallel sequence processing, but standard dense attention has quadratic memory/compute cost in sequence length.
\end{block}
\end{frame}

\begin{frame}{Transformer Families and Use Cases}
\begin{table}
\small
\begin{tabular}{ll}
\toprule
\textbf{Architecture} & \textbf{Typical use} \\
\midrule
Encoder-only & Representation learning / classification \\
Decoder-only & Autoregressive generation \\
Encoder-decoder & Conditional sequence generation \\
Vision transformers & Patch/token-based vision models \\
Multimodal transformers & Joint text/image/audio representations \\
\bottomrule
\end{tabular}
\end{table}

\begin{itemize}
  \item Pretraining and transfer learning are central to many modern Transformer applications.
  \item Long-context performance, latency, memory, data requirements, and task structure should determine whether a Transformer is preferable to simpler alternatives.
\end{itemize}
\end{frame}

% ================================================================
% PART 4: PRACTICAL CONSIDERATIONS
% ================================================================

\begin{frame}{Regularization in Deep Learning}
\begin{columns}
\begin{column}{0.48\textwidth}
\textbf{Weight decay}
\[
  \mathcal{L}_{\text{total}}=\mathcal{L}+\lambda\sum_i w_i^2.
\]
Penalizes large weights and changes the optimization objective.

\vspace{0.5em}
\textbf{Dropout}
\begin{itemize}
  \item randomly masks activations during training;
  \item changes the effective network seen by each minibatch;
  \item dropout probability is a tunable hyperparameter.
\end{itemize}
\end{column}

\begin{column}{0.48\textwidth}
\textbf{Batch normalization}
\begin{itemize}
  \item normalizes intermediate activations using minibatch statistics;
  \item can stabilize optimization and permit different learning-rate regimes;
  \item may have regularizing effects, but regularization is not its sole role.
\end{itemize}

\vspace{0.5em}
\textbf{Data augmentation}
\begin{itemize}
  \item generates plausible transformed examples;
  \item transformations must preserve target semantics.
\end{itemize}
\end{column}
\end{columns}
\end{frame}

\begin{frame}{Early Stopping as Model Selection}
Early stopping selects a checkpoint using validation performance rather than continuing optimization indefinitely.

\begin{figure}
\centering
\begin{tikzpicture}[scale=0.9]
\draw[->] (0,0) -- (6,0) node[right] {Epochs};
\draw[->] (0,0) -- (0,4) node[above] {Loss};
\draw[thick, forest] plot[smooth] coordinates {(0.5,3.5) (1,2.5) (2,1.5) (3,1) (4,0.7) (5,0.5)};
\node at (5,0.2) {Train};
\draw[thick, crimson] plot[smooth] coordinates {(0.5,3.5) (1,2.5) (2,1.7) (3,1.3) (4,1.4) (5,1.6)};
\node at (5,1.9) {Validation};
\draw[dashed] (3,0) -- (3,1.3);
\node at (3,-0.35) {Selected checkpoint};
\end{tikzpicture}
\end{figure}

\begin{alertblock}{Evaluation rule}
The validation set used for early stopping is part of model selection. Final performance should still be estimated on untouched evaluation data.
\end{alertblock}
\end{frame}

\begin{frame}{Knowledge Check: Deep Learning}
\textbf{Test your understanding}

\vspace{0.2cm}

\textbf{Question 1:} Why don't we stack linear layers without activations?

\begin{enumerate}[A)]
\item It's too slow
\item \textcolor{forest}{\textbf{Multiple linear transformations collapse to one linear transformation}}
\item It causes vanishing gradients
\item Backprop doesn't work
\end{enumerate}

\pause

\begin{block}{Explanation}
\textcolor{forest}{\textbf{Answer: B}} - Without non-linear activations, stacking layers is pointless: $f(W_2(W_1 x)) = W_2 W_1 x = Wx$ is still just a linear function. Non-linearity is essential for learning complex patterns. This is why we use ReLU, sigmoid, tanh, etc. between layers.
\end{block}

\pause

\vspace{0.2cm}

\textbf{Question 2:} What problem do ResNets solve with skip connections?

\begin{enumerate}[A)]
\item Overfitting
\item \textcolor{forest}{\textbf{Vanishing gradients in very deep networks}}
\item Slow training
\item Too many parameters
\end{enumerate}

\pause

\begin{block}{Explanation}
\textcolor{forest}{\textbf{Answer: B}} - Skip connections ($y = F(x) + x$) create identity mappings that allow gradients to flow directly through the network during backprop. Without them, gradients diminish exponentially in very deep networks (50+ layers), making training impossible. ResNets enabled 100+ layer networks.
\end{block}
\end{frame}

\begin{frame}{Summary \& Best Practices}
\textbf{Key takeaways for deep learning}

\begin{columns}
\begin{column}{0.5\textwidth}
\textbf{Architecture Choices:}

\begin{itemize}
\item \textbf{Images:} CNNs (ResNet, EfficientNet) or Vision Transformers
\item \textbf{Sequences:} Transformers (BERT, GPT) or LSTMs for small data
\item \textbf{Tabular:} Tree ensembles are strong baselines; neural models should justify added complexity empirically
\item \textbf{Time series:} LSTMs, Transformers, or specialized (N-BEATS)
\end{itemize}

\vspace{0.3cm}

\textbf{Training Best Practices:}

\begin{enumerate}
\item Start simple, add complexity
\item Use pre-trained models when possible
\item Normalize inputs (mean=0, std=1)
\item Treat optimizer and learning-rate schedule as experimental choices; compare plausible baselines under matched budgets
\item Monitor train AND validation loss
\item Use early stopping
\item Batch normalization for stability
\item Dropout for regularization
\item Data augmentation for images
\end{enumerate}
\end{column}

\begin{column}{0.5\textwidth}
\textbf{Common Pitfalls:}

\begin{itemize}
\item[$\times$] Not checking for bugs (use small dataset first)
\item[$\times$] Wrong loss function
\item[$\times$] Learning rate too high/low
\item[$\times$] Not shuffling training data
\item[$\times$] Using test set for validation
\item[$\times$] Forgetting to normalize inputs
\item[$\times$] No baseline comparison
\end{itemize}

\vspace{0.3cm}

\textbf{Resources:}

\begin{itemize}
\item \textbf{Courses:}
  \begin{itemize}
  \item Andrew Ng: Deep Learning Specialization
  \item Fast.ai: Practical Deep Learning
  \item Stanford CS231n (CNNs)
  \item CS224n (NLP with Deep Learning)
  \end{itemize}

\item \textbf{Frameworks:}
  \begin{itemize}
  \item PyTorch (research, flexibility)
  \item TensorFlow/Keras (production)
  \item JAX (performance, functional)
  \end{itemize}

\item \textbf{Papers:}
  \begin{itemize}
  \item "Attention Is All You Need" (Transformers)
  \item "Deep Residual Learning" (ResNet)
  \item Read on arXiv, Papers With Code
  \end{itemize}
\end{itemize}
\end{column}
\end{columns}
\end{frame}

% ================================================================
% END OF DEEP LEARNING PRESENTATION
% ================================================================


